% 彗星（综述）
% license CCBYSA3
% type Wiki

本文根据 CC-BY-SA 协议转载翻译自维基百科\href{https://en.wikipedia.org/wiki/Comet}{相关文章}。

\begin{figure}[ht]
\centering
\includegraphics[width=8cm]{./figures/45ca94c34375617b.png}
\caption{} \label{fig_Comet_1}
\end{figure}
彗星是一类由冰组成的、体积较小的太阳系天体或星际天体。当其经过接近太阳的区域时，会因受热而开始释放气体，这一过程称为释气作用（outgassing）。这一过程会在彗核周围产生一个延展的、未受引力束缚的大气层，即彗发（coma），并有时形成由气体和尘埃组成的彗尾，这些物质自彗发中被吹离。上述现象源于太阳辐射及向外流动的太阳风等离子体对彗核施加的作用。彗核的大小从数百米到数十公里不等，由松散聚集的冰、尘埃以及小型岩石颗粒构成。彗发的尺度可达地球直径的 15 倍，而彗尾的长度甚至可能超过一个天文单位。如果距离足够近且亮度足够高，彗星无需借助望远镜即可从地球上观测到，并可能在天空中呈现出长达 30 度（相当于 60 个满月）的角距。自古以来，彗星已被许多文化和宗教所观察和记录。

彗星通常具有高离心率的椭圆轨道，其轨道周期范围广泛，从数年到可能长达数百万年。短周期彗星来源于柯伊伯带或其相关的散射盘，这些区域位于海王星轨道之外。长周期彗星则被认为起源于奥尔特云，一个由冰质天体组成的球形云层，其范围从柯伊伯带外侧延伸至朝向最近恒星方向一半的距离。长周期彗星通常由经过的恒星引力扰动或银河潮汐触发，开始向太阳方向运动。双曲线轨道彗星可能仅一次穿越内太阳系，然后被抛向星际空间。彗星的出现称为一次显现（apparition）。

多次接近太阳的灭绝彗星（extinct comets）会失去几乎全部挥发性冰与尘埃，并可能逐渐呈现出类似小型小行星的外观。一般认为，小行星与彗星的起源不同：小行星形成于木星轨道以内，而非太阳系外部。然而，主带彗星（main-belt comets）以及活跃的半人马族（centaurs）小天体的发现，使小行星与彗星之间的界限变得模糊。21 世纪初，一些具备长周期彗星轨道、但表现出内太阳系小行星特征的小天体被称为曼岛彗星（Manx comets）。它们仍被归类为彗星，例如 C/2014 S3（PANSTARRS）。在 2013 年至 2017 年间，共发现了 27 颗曼岛彗星。

截至 2021 年 11 月，共已知 4,584 颗彗星。然而，这仅占可能存在的彗星总量的一小部分，因为太阳系外部（奥尔特云）中类彗星天体的储库估计约为一万亿个。平均每年约有一颗彗星可被肉眼看见，尽管其中许多较为暗淡，外观并不特别醒目。特别明亮的彗星被称为“伟大彗星”（great comets）。人类已利用无人探测器访问过彗星，例如美国 NASA 的“深度撞击号”（Deep Impact），其通过撞击坦普尔 1 号彗星形成一个新陨坑以研究彗星内部；又如欧洲航天局（ESA）的“罗塞塔号”（Rosetta），其成为首个在彗星上着陆的机器人探测器。
\subsection{词源}
\begin{figure}[ht]
\centering
\includegraphics[width=8cm]{./figures/6a3173d4cf8911ea.png}
\caption{公元 729 年的《盎格鲁－撒克逊编年史》以及尊敬的比德（the Venerable Bede）的著作中均提到了彗星。} \label{fig_Comet_2}
\end{figure}
单词“comet（彗星）”源自古英语 cometa，其来源为拉丁语 comēta 或 comētēs，而这些拉丁词又是对希腊语 κομήτης 的罗马化形式，意为“留着长发的（事物）”。《牛津英语词典》指出，在希腊语中术语（ἀστὴρ） κομήτης 已经表示“长发之星，即彗星”。κομήτης 来源于动词 κομᾶν（koman），意为“留长发”，而该动词又派生自名词 κόμη（komē），意为“头发”，此词也被用来表示“彗尾”。

彗星的天文学符号（在 Unicode 中表示）为 U+2604 ☄ COMET，由一个小圆盘与三条类似头发的延伸构成。
\subsection{物理特征}
\begin{figure}[ht]
\centering
\includegraphics[width=8cm]{./figures/37fee7dd52548b31.png}
\caption{} \label{fig_Comet_3}
\end{figure}